\documentclass{article}

\usepackage{float}
\usepackage{graphicx}
\usepackage{amsmath}
\usepackage{amssymb}
\usepackage{booktabs}
\usepackage{xcolor}
\usepackage{hyperref}
\usepackage{listings}
\usepackage{enumitem}
\usepackage[backend=bibtex,style=numeric]{biblatex}
\addbibresource{bib.bib}

\definecolor{codegreen}{rgb}{0,0.6,0}
\definecolor{codegray}{rgb}{0.5,0.5,0.5}
\definecolor{codepurple}{rgb}{0.58,0,0.82}
\definecolor{backcolour}{rgb}{0.95,0.95,0.92}

\lstdefinestyle{mystyle}{
	backgroundcolor=\color{backcolour},   
	commentstyle=\color{codegreen},
	keywordstyle=\color{magenta},
	numberstyle=\tiny\color{codegray},
	stringstyle=\color{codepurple},
	basicstyle=\ttfamily\footnotesize,
	breakatwhitespace=false,         
	breaklines=true,                 
	captionpos=b,                    
	keepspaces=true,                 
	numbers=left,                    
	numbersep=5pt,                  
	showspaces=false,                
	showstringspaces=false,
	showtabs=false,                  
	tabsize=2
}

\lstset{style=mystyle}

\title{\textbf{Political Drift Model: Revised Formalization and Implementation}}
\author{Luis Luarte}
\date{\today}

\begin{document}
	
	\maketitle
	
	\begin{abstract}
		This document presents a revised mathematical formalization of the Political Drift Model, which characterizes legislative behavior as a dynamic trajectory through policy space, and proposes an empirical test. I introduce a theoretical framework in which agent positions evolve through a dual mechanism: an adaptive learning process based on voting outcomes and a "force field" representing institutional party influence. This revision focuses on the precise mathematical definition of these governing dynamics and the structural constraints required for robust parameter estimation.
	\end{abstract}
	
	\section{Introduction}
	Understanding the evolution of legislative behavior requires analyzing how political agents adjust their ideological positions over time. This report details a dynamic model of "Political Drift," in which a legislator's position is not static but evolves along a trajectory through a policy space.
	
	The model assumes that two primary factors influence an agent's state at any given time:
	\begin{enumerate}
		\item \textbf{Learning from Outcomes:} Agents adjust their position based on the success or failure of their votes relative to the final legislative outcome.
		\item \textbf{Institutional Gravity (Force Field):} Agents are subject to "gravitational" pulls from political party centers. This force field acts as a constraining or attracting mechanism, representing party discipline or ideological alignment.
	\end{enumerate}
	
	By formalizing these dynamics, I aim to recover the latent parameters that govern susceptibility to party influence and the rate of ideological adaptation using Maximum Likelihood Estimation (MLE).
	
	\section{Methods}
	
	\subsection{Data and Sampling Strategy}
	The dataset consists of legislative voting records organized by period. For each politician and period, their political party was predefined, and political ideology was estimated using dynamic ideal-point estimation (this data was recovered from the course repository). For each party, a mean political center was computed by averaging over all politicians in that party at each legislative voting period. To ensure the computational feasibility of dynamic estimation while preserving the temporal topology of the data, I employ a stratified hierarchical sampling method. I sampled $N = 10$ politicians per period and considered the first $1000$ bills.
	
	\paragraph{Political Roles}
	To analyze the behavior of different political cohorts, agents are classified into two roles based on their relationship to the executive branch at time $t$. An agent is defined as \textbf{Officialist} if their party affiliation matches the party of the sitting President (the ruling party). An agent is defined as \textbf{Opposition} if they belong to any party other than the ruling party. This classification is dynamic and updates if the ruling party changes between periods.
	
	\subsection{Mathematical Model}
	
	The core of the model is the evolution of an agent's ideal point $p_{i,t}$. Voting decisions drive this evolution at each step, and are influenced by external pressures from the party system.
	
	\subsubsection{Voting Behavior and Probability}
	At every time step $t$, an agent $i$ faces a vote choice between a proposal $x_p$ and a status quo $x_{sq}$. 
	
	The spatial locations of these parameters are estimated based on the assumption of partisan agenda control. The proposal $x_p$ is modeled as the ruling party's mean ideological position, reflecting the executive's agenda-setting power to introduce bills aligned with the coalition's preferences. The status quo $x_{sq}$ is modeled as the mean position of the entire legislature, serving as a proxy for the pre-existing legislative consensus or equilibrium.
	
	The decision to vote "Yea" or "Nay" is determined by a utility function that weighs ideological proximity against party loyalty (portfolio payoffs).
	
	The total utility for voting "Yea" ($U_{i,\text{yea}}$) or "Nay" ($U_{i,\text{nay}}$) is defined as a convex combination of ideological and portfolio utilities, weighted by a parameter $\omega \in [0,1]$:
	
	\begin{equation}
		U_{i, \text{choice}} = \omega \cdot u^{\text{ideo}}_{i, \text{choice}} + (1 - \omega) \cdot u^{\text{port}}_{i, \text{choice}}
	\end{equation}
	
	\textbf{Ideological Utility:} This is modeled as a quadratic loss function based on the distance between the agent's position $p_{i,t}$ and the policy location:
	\begin{align}
		u^{\text{ideo}}_{i, \text{yea}} &= -(p_{i,t} - x_p)^2 \\
		u^{\text{ideo}}_{i, \text{nay}} &= -(p_{i,t} - x_{sq})^2
	\end{align}
	
	\textbf{Portfolio Utility:} This represents party discipline. Let $S_{i,t} \in \{\text{officialist}, \text{opposition}\}$ be the agent's role relative to the ruling party. The party loyalty parameter $K$ confers a benefit for voting with one's coalition and a penalty for voting against it:
	\begin{equation}
		u^{\text{port}}_{i, \text{yea}} = 
		\begin{cases} 
			K & \text{if } S_{i,t} = \text{officialist} \\
			-K & \text{if } S_{i,t} = \text{opposition}
		\end{cases}
	\end{equation}
	
	The probability of voting "Yea" is then derived using the logistic function of the utility difference:
	\begin{equation}
		P(\text{vote}_i = \text{Yea}) = \frac{1}{1 + \exp\left(-(U_{i,\text{yea}} - U_{i,\text{nay}})\right)}
	\end{equation}
	
	\subsubsection{Dynamic Trajectory}
	The agent $i$'s ideological position at time $t$, denoted as $p_{i,t}$, evolves according to a learning component and a force field component:
	\begin{equation}
		p_{i, t+1} = p_{i,t} + \Delta_{\text{learn}} + \Delta_{\text{force}}
	\end{equation}
	
	\paragraph{Learning Component ($\Delta_{\text{learn}}$)}
	Agents update their beliefs based on whether their vote aligned with the legislative outcome. Let $x_{\text{winner}}$ be the policy location of the winning option ($x_p$ or $x_{sq}$). The learning update is proportional to the distance between the agent's current position and the winning policy, modulated by a learning rate $\alpha$:
	
	\begin{equation}
		\Delta_{\text{learn}} = \alpha \cdot (x_{\text{winner}} - p_{i,t})
	\end{equation}
	
	The learning rate $\alpha$ depends on the agent's success: $\alpha_{\text{win}}$ if the agent voted for the winner, and $\alpha_{\text{loss}}$ otherwise. This allows for asymmetric learning dynamics (e.g., agents may react more strongly to losing).
	
	\subsubsection{The Force Field}
	The force exerted on an agent by the set of party centers $C_t = \{C_1, \dots, C_M\}$ in period $t$ is calculated as a summation of inverse-square forces. Let $\theta$ represent the mass/gravity of the parties and $\kappa$ represent the agent's susceptibility.
	\begin{equation}
		\Delta_{\text{force}} = \kappa \sum_{j \in C_t} \frac{(C_j - p_{i,t}) \cdot \theta^\eta}{(|C_j - p_{i,t}| + \epsilon)^3}
	\end{equation}
	Where $\epsilon = 0.001$ is a smoothing parameter to prevent singularities, and $\eta$ controls the scaling of the mass (assumed $\eta=2$ in implementation). This is a slight change from initial model specification, previously I considered $\theta_{party}$ and $\theta_{others}$ to represent that ``my'' party could exert a stronger pull, however, I figured this is taken care of by the ``most likely political definition''  $j_{\text{self}}(t) \gets \underset{j \in J}{\text{argmin}} |C_j(t) - p_i(t)|$. Also, $\eta$ was set to a constant rather than a free parameter, as $k$ already captures the susceptibility to the political space and potential ``pull'' from the parties. This also prevents a parameter identifiability issue during optimization.
	
	\subsection{Likelihood Function}
	The model parameters $\Theta = \{\omega, K, \alpha_{\text{win}}, \alpha_{\text{loss}}, \kappa, \theta\}$ are estimated by minimizing the negative log-likelihood of the observed voting sequence. For a sequence of $T$ votes $v_1, \dots, v_T$:
	
	\begin{equation}
		\mathcal{L}(\Theta) = - \sum_{t=1}^{T} \ln P(\text{vote}_i = v_t | p_{i,t}, \Theta)
	\end{equation}
	
	The trajectory $p_{i,t}$ is recursively computed at each step based on the previous state and the parameters, making the likelihood computation dynamic.
	
	\subsection{Computational Optimization}
	
	For parameter optimization, I used R's `nlminb`, a Quasi-Newton method that natively handles box constraints. This ensures that parameters such as weights ($w \in [0,1]$) and susceptibility ($\kappa \geq 0$) remain within theoretically valid bounds without the need for manual penalty functions.
	
	\section{Results}
	
	\subsection{Hypothesis 1: Asymmetric Learning}
	To determine whether legislative agents exhibit asymmetric learning dynamics, we compared the estimated learning rates for winning ($\alpha_{\text{win}}$) and losing ($\alpha_{\text{loss}}$) votes. We conducted a Wilcoxon signed-rank test with continuity correction to test the null hypothesis that the actual location shift between the paired parameters is zero.
	
	The analysis revealed a statistically significant difference between the learning rates ($V = 45$, $p = 0.02633$). This result allows us to reject the null hypothesis at the $\alpha = 0.05$ significance level. The test statistic $V = 45$ (based on the sample size $N=10$) indicates a high sum of positive ranks, suggesting that for the majority of sampled agents, one learning rate (typically $\alpha_{\text{loss}}$ in this context) consistently exceeds the other.
	
	\subsection{Hypothesis 2: Systemic Drift}
	To test for systemic drift relative to the governing coalition, we employed a hierarchical linear mixed-effects model (LME). This approach allows us to estimate drift trajectories relative to the \textbf{Ruling Party Political Center ($x_p$)} while accounting for the nested structure of the data (observations of $p_{it}$ nested within agents) and separating individual variation (random effects) from systemic role-based trends (fixed effects).
	
	The model specification was:
	\begin{equation}
		|p_{it} - x_{p,t}| = \beta_0 + \beta_1 \ln(t) + \beta_2 \text{Role}_i + \beta_3 (\ln(t) \times \text{Role}_i) + u_{0i} + u_{1i} \ln(t) + \epsilon_{it}
	\end{equation}
	where $|p_{it} - x_{p,t}|$ represents the absolute ideological distance from the ruling party, and $\ln(t)$ captures potential effects of passing periods (diminishing returns of learning).
	
	The analysis of Estimated Marginal Trends revealed a significant divergence for opposition agents. While the officialist group showed a mean drift slope indistinguishable from zero ($Slope = 0.008, p = 0.63$), indicating static behavior relative to its own party (officialist is defined as being within the ruling party), the opposition group exhibited a statistically significant positive slope ($Slope = 0.09, p < 0.01$). The magnitude and direction clearly point towards divergence: opposition agents are drifting \textit{away} from the ruling party's political center over time.
	
	\subsection{Hypothesis 3: Agent Typologies}
	
	To investigate the underlying incentives driving these behaviors, we analyzed the joint distribution of the estimated parameters for ideological weight ($\omega$) and party loyalty ($K$). The results reveal a structural separation between the two groups (Figure \ref{fig:p1}). 
	
	Officialists (red) cluster in the top-left quadrant, characterized by high party loyalty ($K > 0.5$) and lower ideological weights. This suggests their behavior is dominated by portfolio incentives, voting with the coalition to maximize loyalty payoffs. In contrast, Opposition agents (blue) cluster in the bottom-right quadrant, exhibiting high ideological weights ($\omega > 0.5$) and low or negative loyalty parameters. This implies they act as ``ideologues'' unconstrained by party discipline and primarily motivated by spatial policy maximization. Statistical tests showed that $K_{\text{off - opp}}$, and $\omega_{\text{off - opp}}$ were statistically significant ($p < 0.01$). These results suggest that the partition shown graphically is also statistically significant.
	
	\begin{figure}[H]
		\centering
		\includegraphics[width=\textwidth]{p1.pdf}
		\caption{Joint distribution of estimated Ideological Weight ($\omega$) and Party Loyalty ($K$) parameters by political role. The plot reveals distinct agent typologies: Officialists (red) behave as disciplined party members (High $K$, Low $\omega$), while Opposition agents (blue) behave as unconstrained ideologues (High $\omega$, Low $K$).}
		\label{fig:p1}
	\end{figure}
	
	\section{Discussion}
	The detection of asymmetric learning rates ($\alpha_{\text{loss}} > \alpha_{\text{win}}$) supports the theoretical proposition of negativity bias: legislators are more reactive to policy failures. However, contrary to the initial hypothesis that this learning mechanism would drive convergence towards the powerful agenda-setter (the Ruling Party), the systemic drift analysis revealed a significant \textit{divergence} for Opposition agents. They drift away from the ruling party over time, while Officialists remain static.
	
	This divergence is explained by the distinct agent typologies uncovered in Hypothesis 3. The structural separation in parameters ($K$, $\omega$) confirms that agents are rational actors maximizing utility within a complex field of forces. Officialists behave as disciplined "Party Soldiers": their high loyalty payoff ($K$) and low ideological weight ($\omega$) anchor them to the party line, rendering their spatial position ($p_i$) irrelevant to their utility and thus preventing drift. Conversely, Opposition agents behave as "Ideologues": lacking the loyalty incentives ($K$) of the coalition, they maximize utility purely through spatial positioning ($\omega \approx 1$). 
	
	If agents were merely generic representatives, we would expect similar $K$ and $\omega$ values across roles. Instead, the observed divergence is a rational response to the institutional context. For the Opposition, drifting away from a centrist Ruling Party reduces the utility cost of voting "Nay," allowing them to "learn to oppose" effectively. This coherence between the drift trajectories and the underlying parameter estimates validates the model's specification of the political field as a dual system of institutional gravity and adaptive learning.
	
	The limitations of this work can be summarised mainly as computational constraints. Sampling methods were used to estimate parameters; instead of a Bayesian-like optimization that would allow me to derive uncertainty from the parameters, I opted for faster optimization using `nlminb`. I also considered all the bills to be from the ruling party. I determined the voting rule as simple majority in all cases, as I reasoned that agents observe voting proportions to determine were the ``scale'' is leaning towards, and as a first approach to a more agent-based modeling principle were the agent must derive through experience what is observing, rather than being delivered noisyless information, as it would be deriving each bill status quo or political position through semantic and contextual analysis, and delivered equally to all agents. Finally, much work is to be done to validate model specifications, (1) no parameter sensitivity was performed, (2) no competing ``null'' model was set as reference.
	
	
\end{document}